\section{Requerimientos candidatos}

La primera lectura del caso permite identificar siete necesidades de alto nivel
que posteriormente se precisan como requerimientos funcionales y no
funcionales:

\begin{enumerate}[label=\textbf{RC-\arabic*.}, leftmargin=*]
    \item Almacenar información de sucursales, premios y clientes.
    \item Agregar, consultar, editar y eliminar registros de las tres entidades.
    \item Conservar la información entre ejecuciones mediante archivos CSV.
    \item Proporcionar al menos un menú de interacción para ejecutar las
    operaciones solicitadas.
    \item Recibir la llave de una entidad y devolver todos los datos asociados
    con el registro.
    \item Aceptar únicamente valores numéricos en los campos que así lo
    requieran.
    \item Manejar excepciones e informar al usuario cuando una operación no
    pueda completarse.
\end{enumerate}

Estos candidatos delimitan el objetivo inmediato del prototipo. Aspectos como
la integridad de los identificadores, el formato de los demás atributos, la
recuperación ante errores y la organización interna del código se detallan en
las secciones siguientes.
